\chapter{Cohomology and Explicit Deformations}

Work over a field \(k\) of characteristic zero. For Lie brackets,
our convention for the Chevalley--Eilenberg differential, changes
of basis, and obstruction classes is that of
\emph{Lie Algebras and Finite-Dimensional Representations},
Results I02--I03~\cite[PDF pp.~56--57]{mg-lie}.
We use those results to compute an entire first-order obstruction
locus in a small example. We then turn to Lie bialgebras and
deformations of the coproduct. In that setting multiplication
and comultiplication must satisfy a compatibility equation
in addition to their separate associativity conditions.

\section{An adjoint cohomology calculation}

\begin{proposition}[Adjoint cohomology of the three-dimensional Heisenberg algebra]
\label{res:LD01}
Let \(\mathfrak h\) have basis \(x,y,z\) with
\([x,y]=z\) and \(z\) central. Then
\(\dim H^2(\mathfrak h,\mathfrak h)=5\) and
\(\dim H^3(\mathfrak h,\mathfrak h)=2\), with adjoint
coefficients. Each class in degree two has a unique representative
\(\phi\) of the form
\[
 \phi(x,y)=0,\qquad
 \phi(x,z)=a x+b y+c z,\qquad
 \phi(y,z)=d x-a y+e z .
\]
In degree three evaluation on \((x,y,z)\) identifies
\(H^3\) with \(\mathfrak h/kz\).
\end{proposition}
\begin{proof}
For an alternating \(2\)-cochain \(\psi\), write
\(A=\psi(x,y)\), \(B=\psi(x,z)\), and \(C=\psi(y,z)\).
The argument with a bracket in the Chevalley--Eilenberg
formula is zero: the only nonzero bracket of two basis vectors
in the triple is \([x,y]=z\), and \(\psi(z,z)=0\).
Hence
\[
 (d\psi)(x,y,z)=[x,C]-[y,B]
               =(C_y+B_x)z .
\]
Here subscripts denote coordinates in the fixed basis.
Thus the cocycle space has dimension \(8\). The image in
degree three is exactly \(kz\), and there are no degree-four
cochains, proving the degree-three assertion.

For a linear \(1\)-cochain \(f\) write
\(f(z)=u x+v y+w z\). Its differential is
\[
 \begin{aligned}
 (df)(x,y)&=-u x-v y+(f(y)_y+f(x)_x-w)z,\\
 (df)(x,z)&=v z,\\
 (df)(y,z)&=-u z.
 \end{aligned}
\]
Thus a coboundary is determined by its value on \((x,y)\),
and that value ranges over all of \(\mathfrak h\), for instance
by taking \(f(x)=f(y)=0\) and varying \(f(z)\).
Subtract the unique coboundary matching \(A\). In the resulting
normalized cocycle \(A=0\), the only equation is
\(C_y=-B_x\). This gives the claimed parameters and uniqueness.
\end{proof}

\begin{theorem}[Exactly which of these first-order classes extend]
\label{res:LD02}
In the coordinates of Proposition~\ref{res:LD01}, a first-order
bracket deformation of \(\mathfrak h\) extends modulo \(t^3\)
if and only if
\[
                  cd-ea=0,\qquad eb+ca=0 .
\]
When both equations hold, the bracket \([\, ,\,]+t\phi\)
already satisfies Jacobi over \(k[t]\), and in particular gives
a formal extension over \(k[[t]]\).
\end{theorem}
\begin{proof}
First take the unique normalized representative. In Jacobi,
use the cyclic expression \(\sum_{\mathrm{cyc}}[[x,y],z]\)
as in Results I02--I03 of \cite{mg-lie}.
The coefficient in order \(t\) is zero since \(\phi\) is
a cocycle. At order \(t^2\), the contribution of the first-order
cochain on \((x,y,z)\) is
\[
 \begin{aligned}
   \sum_{\mathrm{cyc}}\phi(\phi(x,y),z)
    &= -e\,\phi(x,z)+c\,\phi(y,z)\\
    &= (cd-ea)x-(eb+ca)y .
 \end{aligned}
\]
If a \(t^2\)-cochain is added to the bracket, its contribution
at that order is the negative of its differential, by
Result I03. By Proposition~\ref{res:LD01}, every such differential
takes values in \(kz\) on \((x,y,z)\). It can cancel the
displayed vector only if both its coordinates are zero.

When both vanish, Jacobi for the bracket with only the linear
term is identically zero: its coefficients have degrees \(0,1,2\),
and each has vanished on an exterior basis for triples.
Changing the representative of the first-order class by a
coboundary preserves existence of either extension. The first-order
change is \(1+t f\), as in Result I02, and the same linear
operator is invertible modulo \(t^3\) and over \(k[[t]]\) by
inverting it as a formal power series. Transport the extension
by that change of basis.
\end{proof}

\begin{example}[Two extendable classes with an obstructed sum]
\label{res:LD03}
Let \(\phi_1(x,z)=x,\phi_1(y,z)=-y\), and
\(\phi_2(x,z)=0,\phi_2(y,z)=z\), with both cochains zero
on \((x,y)\). Each extends, but their sum is obstructed
already in degree two.
The extension defined by \(\phi_1\) specializes at \(t=0\)
to \(\mathfrak h\) and is isomorphic to
\(\mathfrak{sl}_2(k((t)))\) over the Laurent series field.
\end{example}
\begin{proof}
For \(\phi_1\) only \(a=1\) is nonzero, and for \(\phi_2\)
only \(e=1\) is nonzero. The equations of
Theorem~\ref{res:LD02} both hold in either case.
For their sum the obstruction vector is \(-x\), nonzero
in the quotient in Proposition~\ref{res:LD01}.
For the generic \(\phi_1\)-bracket put
\(H=-2z/t\), \(E=x\), \(F=-2y/t\). These form a basis over
\(k((t))\), and its brackets are \([H,E]=2E\),
\([H,F]=-2F\), and \([E,F]=H\), precisely the displayed
matrix-basis brackets of \(\mathfrak{sl}_2\).
\end{proof}

The general obstruction formalism for deformations over Artinian
bases is developed by Fialowski--Fuchs~\cite[Sections 1 and 3]{fialowski-fuchs}.
The calculation above concerns extensions of each first-order class
over one parameter.

\section{Coboundary Lie bialgebras}

\begin{definition}[Lie bialgebras and coboundary cobrackets]
\label{res:LD04}
For a finite-dimensional Lie algebra \(\mathfrak g\), put
\(x\cdot(y\wedge z)=[x,y]\wedge z+y\wedge[x,z]\).
A Lie bialgebra cobracket is a linear map
\(\delta:\mathfrak g\to\bigwedge^2\mathfrak g\)
such that
\[
 \delta([x,y])=x\cdot\delta(y)-y\cdot\delta(x)
\]
and the dual operation, given by
\(\langle[\alpha,\beta]_\delta,x\rangle
=\langle\alpha\wedge\beta,\delta(x)\rangle\), is a Lie bracket.
For \(r\in\bigwedge^2\mathfrak g\) the coboundary map
\(\delta_r(x)=x\cdot r\) always satisfies the first condition.
\end{definition}

\begin{proposition}[The invariant Schouten-square criterion]
\label{res:LD05}
Define a symmetric bracket on bivectors by bilinear extension of
\[
 \begin{aligned}
 [a\wedge b,c\wedge d]_{\mathrm S}
   ={}&[a,c]\wedge b\wedge d-[a,d]\wedge b\wedge c\\
      &-[b,c]\wedge a\wedge d+[b,d]\wedge a\wedge c .
 \end{aligned}
\]
The map \(\delta_r\) is a Lie bialgebra cobracket if and only if
\([r,r]_{\mathrm S}\) is fixed by the adjoint action
on \(\bigwedge^3\mathfrak g\).
\end{proposition}
\begin{proof}
The displayed formula is alternating in the two factors of
each bivector, so it descends to the exterior squares.
Interchanging the two bivectors changes the sign of the Lie
bracket and of the necessary permutation of the two remaining
vectors, making the formula symmetric. Applying
\(\operatorname{ad}_x\) to each term gives
\[
 x\cdot[r,s]_{\mathrm S}=[x\cdot r,s]_{\mathrm S}
                         +[r,x\cdot s]_{\mathrm S},
\]
using the Lie Jacobi identity in the bracketed factor.

For any \(\delta\), extend it in degree two by
\(\partial_\delta(y\wedge z)=\delta(y)\wedge z-y\wedge\delta(z)\).
Pairing \(\partial_\delta\delta(x)\) with \(\alpha\wedge\beta\wedge\gamma\)
gives the pairing of \(x\) with the cyclic sum
\(\sum_{\mathrm{cyc}}[[\alpha,\beta]_\delta,\gamma]_\delta\).
Indeed for a term \(u\wedge v\) of \(\delta(x)\) the pairing of
\(\delta(u)\wedge v-u\wedge\delta(v)\) is the cyclic sum of
\(
\langle[\alpha,\beta]_\delta,u\rangle\langle\gamma,v\rangle
-\langle[\alpha,\beta]_\delta,v\rangle\langle\gamma,u\rangle
\).
These pairings detect every trivector, so the dual Jacobi
identity is equivalent to \(\partial_\delta\delta=0\).

The bivector formula gives
\([r,y\wedge z]_{\mathrm S}=-\delta_r(y)\wedge z+
y\wedge\delta_r(z)\), first for a decomposable \(r\)
and then by linearity. Consequently
\[
 \partial_{\delta_r}\delta_r(x)
      =-[r,x\cdot r]_{\mathrm S}
      =-\tfrac12 x\cdot[r,r]_{\mathrm S}.
\]
The \(1\)-cocycle condition also holds because the adjoint
action is a representation. This proves the criterion.
\end{proof}

\begin{example}[Two cobrackets on the split rank-one algebra]
\label{res:LD06}
For the basis \(H,E,F\) with \([H,E]=2E\),
\([H,F]=-2F\), \([E,F]=H\), there are Lie bialgebra
structures from \(r_1=F\wedge E\) and \(r_2=H\wedge E\).
The first has
\(\delta_{r_1}(H)=0\), \(\delta_{r_1}(E)=H\wedge E\),
and \(\delta_{r_1}(F)=H\wedge F\).
The second has Schouten square zero and has
\(\delta_{r_2}(H)=2H\wedge E\), \(\delta_{r_2}(E)=0\),
and \(\delta_{r_2}(F)=-2E\wedge F\).
\end{example}
\begin{proof}
For \(r=u\wedge v\), the formula of
Proposition~\ref{res:LD05} gives
\([r,r]_{\mathrm S}=2[u,v]\wedge u\wedge v\).
For \(r_1\), this is a scalar multiple of \(H\wedge E\wedge F\).
The adjoint action of \(H\) on that trivector cancels the
weights \(2,-2\). The adjoint actions of \(E\) and \(F\)
each give repeated factors, so the trivector is invariant.
For \(r_2\), \([H,E]=2E\) gives a repeated factor in the square.
The criterion gives both bialgebras. Acting on the two
bivectors with the basis vectors gives the claimed formulas.
The first convention agrees with the simple-root normalization
of the preceding quantum chapter.
\end{proof}

\section{First-order changes of a bialgebra}

\begin{proposition}[Linearized associativity, coassociativity, and compatibility]
\label{res:LD07}
Let \(B\) be a bialgebra with multiplication \(m\), coproduct \(\Delta\),
unit \(1\), and counit \(\epsilon_B\).
Over \(k[s]/(s^2)\), write \(m_s=m+sf\), \(\Delta_s=\Delta+sg\)
and keep the unit and counit fixed. Let
\(J=1\otimes\tau\otimes1\) on \(B^{\otimes4}\), where
\(\tau\) interchanges the middle factors.
The bialgebra identities hold precisely when
\[
 \begin{aligned}
 &f(1,b)=f(b,1)=0,\quad \epsilon_B f=0,\quad
   g(1)=0,\quad (\epsilon_B\otimes1)g=(1\otimes\epsilon_B)g=0,\\
 &a f(b,c)-f(ab,c)+f(a,bc)-f(a,b)c=0,\\
 &(\Delta\otimes1)g+(g\otimes1)\Delta
        -(1\otimes\Delta)g-(1\otimes g)\Delta=0,\\
 &\Delta f+gm=(f\otimes m+m\otimes f)J(\Delta\otimes\Delta)
             +(m\otimes m)J(g\otimes\Delta+\Delta\otimes g).
 \end{aligned}
\]
Under pullback by \(1+s u\) with \(u(1)=0\) and
\(\epsilon_B u=0\), the two cochains change by
\[
 \begin{aligned}
 f&\longmapsto f+m(u\otimes1+1\otimes u)-u m,\\
 g&\longmapsto g+\Delta u-(u\otimes1+1\otimes u)\Delta.
 \end{aligned}
\]
\end{proposition}
\begin{proof}
Keeping the unit and counit supplies the first line, including
the condition that the coproduct is unital and the counit is
multiplicative. For associativity compare
\(m_s(m_s\otimes1)\) and \(m_s(1\otimes m_s)\);
their difference in degree one gives the second line up to
an overall sign. The degree-one part of
\((\Delta_s\otimes1)\Delta_s=(1\otimes\Delta_s)\Delta_s\)
is the third line. For compatibility, start from
\(\Delta_s m_s=(m_s\otimes m_s)J(\Delta_s\otimes\Delta_s)\).
On the left, the \(s\)-terms are \(\Delta f\) and \(gm\).
On the right one changes one of the four indicated operations
at a time. This gives the last line, and terms with two changes
are zero. These expansions exhaust the bialgebra identities.
In the pullback formulas use the inverse \(1-su\). Apply it
after multiplication and in both factors after comultiplication;
their linear terms are exactly the two displayed changes.
\end{proof}

These equations give the two operation components and change of
coordinates of the normalized bialgebra deformation complex.
For the general cohomology and higher brackets see
Markl~\cite[Section 6]{markl}.

\begin{theorem}[A nontrivial coproduct deformation of a polynomial Hopf algebra]
\label{res:LD08}
Let \(B=k[x,y]\) with \(x,y\) primitive. With multiplication
unchanged, the \(t\)-adically completed algebra \(B[[t]]\)
admits a Hopf structure extending the original one:
\[
 \Delta_t(x)=x\otimes1+1\otimes x,\qquad
 \Delta_t(y)=y\otimes1+\exp(tx)\otimes y .
\]
The counit still vanishes on \(x,y\). Its first-order class
as a bialgebra deformation is nonzero. On primitives of
degree one the linearized skew coproduct gives
\(\delta(x)=0\) and \(\delta(y)=x\wedge y\).
\end{theorem}
\begin{proof}
The formulas determine a continuous algebra map to
\((B\otimes B)[[t]]\), since the target is commutative and
both images have prescribed constant terms. As \(x\) is
primitive, the formal series \(a=\exp(tx)\) is group-like:
the exponential of the sum of the two commuting copies of
\(tx\) is the product of their exponentials.
Both iterated coproducts of \(y\) are
\[
 y\otimes1\otimes1+a\otimes y\otimes1+a\otimes a\otimes y,
\]
and they agree on \(x\) by primitivity. The counit
identities follow on generators. Define the continuous algebra
map \(S\) by \(S(x)=-x\), \(S(y)=-a^{-1}y\), so \(S(a)=a^{-1}\).
The two antipode identities hold on \(x\) and \(y\):
on \(y\) the sums are \(S(y)+a^{-1}y=0\) and
\(y+aS(y)=0\). The two sides are algebra maps by commutativity,
so the identities hold throughout the completion.

At first order the coefficient of \(t\) sends \(y\) to
\(x\otimes y\). The difference of this tensor and its flip
is \(x\otimes y-y\otimes x\), identified with \(x\wedge y\).
Since
\(\Delta\) is cocommutative, the second change of cochains
in Proposition~\ref{res:LD07} is symmetric under the flip,
for every \(u\). Thus the skew part cannot be removed by
a first-order change of coordinates. Its values on \(x,y\)
give the asserted cobracket on the abelian two-dimensional
Lie algebra they span. Its dual is a Lie bracket, since the
Jacobiator is alternating on a two-dimensional space.
\end{proof}

For the quantization problem with a prescribed Lie bialgebra
and its classical limit, see the definition and existence
theorem of Etingof--Kazhdan
~\cite[Section 3 and Theorem 9.6]{etingof-kazhdan}.
